\chapter{Modelo del robot Pioneer}

El Pioneer P3DX se modela como robot diferencial. Su movimiento se controla mediante dos ruedas motrices: izquierda y derecha. A partir de una velocidad lineal deseada \(v\) y una velocidad angular deseada \(\omega\), las velocidades de rueda pueden aproximarse como:

\[
v_L = v - \frac{L}{2}\omega,\qquad
v_R = v + \frac{L}{2}\omega
\]

donde \(L\) es la distancia entre ruedas. En CoppeliaSim estas velocidades se transforman a velocidad angular de motor dividiendo por el radio de rueda \(r\):

\[
\omega_L = \frac{v_L}{r},\qquad \omega_R = \frac{v_R}{r}
\]

\section{Sensores}

Para la segunda parte se emplean sensores de proximidad frontales o laterales. Cada lectura de obstaculo se transforma en una componente repulsiva. La suma de atraccion al objetivo y repulsion de obstaculos produce un vector de movimiento reactivo. Este enfoque es sencillo y adecuado para la actividad, aunque puede presentar minimos locales, oscilaciones o trayectorias suboptimas \cite{siegwart2011,tzafestas2014}.
